% =========================================================================
\chapter{Experimental Design}
\label{ch:design}
% =========================================================================

\claim{Power arithmetic precedes endpoints: on a bed whose same-config flip
rate is $20\%$, total-score deltas of five tasks or fewer are unreadable
--- so the design pairs, pre-registers, interleaves windows, counts
firings, and audits the evaluator itself.}

Most experimental sections describe what was measured. This one begins by
establishing what \emph{could} have been measured, because on this bed that
question has an uncomfortable answer and every subsequent design decision
follows from it.

% -------------------------------------------------------------------------
\section{Beds, arms, and the switch ladder}
\label{sec:beds}
% -------------------------------------------------------------------------

The bed is the \textbf{100-task no-pixel subset} of GAIA's text-only
split; \S\ref{sec:what-was-run} gives the four reasons for that benchmark
and those tasks, of which the one that matters here is headroom --- 97 of
the 100 have been solved at least once against a best single round of 76
\F{16}. The comparison is
between the official system with every GHX flag off (the \emph{no-graph}
arm) and the same system with the GHX switches enabled level by level ---
recording, then graph evidence, then graph verification, then the typed
candidate surface --- each level keeping the ones below it, never all at
once.

Six campaigns are claim-bearing: the no-graph baseline, the post-fix graph
campaign, and two replication pairs re-flying both arm configurations
unmodified for sixteen rounds each \F{44}\F{45}. Two earlier graph
campaigns, flown before the graph layer was correct, are excluded from
every number and every narrative by a scope ruling (\S\ref{sec:scope}).
Scoring is audited throughout.

\subsection*{One denominator for everything}

The arms did not fly quite the same bed. GAIA's text-only split holds 103
tasks, three of which carry pixel content the graph recorder could not
attribute; the no-graph arm attempted all 103 on the first two seeds, the
graph arms attempted the 100 without them, and the third seed's no-graph
arm was moved onto the 100 at round 7, mid-campaign. The loop's own
\texttt{curves.json} records whichever bed each round actually flew, so
\emph{those numbers are not comparable across arms} --- and on the third
seed not even with themselves across the switch.

Every cross-arm reading in this thesis is therefore recomputed from
\texttt{task\_history} on the \textbf{100-task no-pixel subset}, the
ordered subset common to all six campaigns \F{46}.

Of those 100 tasks, \textbf{97 are reachable}: three are never solved in
9{,}600 task evaluations across the six campaigns and contribute no
variance at all, so every \emph{difference} reported in this thesis is
identical whether computed on 97 or on 100, and only the percentages would
shift \F{54}. The 100-task denominator is kept anyway, because dropping
tasks on the basis of their outcomes is selection, however harmless it
looks.

\subsection*{Models, and why this pair}

Every campaign runs the same two models through a LiteLLM gateway, and the
runbooks record the pair on every launch line: the \textbf{solver} is
\texttt{deepseek-v4-flash} (284B parameters, 13B active; thinking enabled,
effort high), and the \textbf{four meta-roles and the judge} are
\texttt{deepseek-v4-pro} (1.6T, 49B active). The official paper's one
complete experiment has the same shape --- a stronger meta tier over a
weaker task tier, Claude Opus over Sonnet, GPT and Qwen --- on a different
benchmark, so the comparison across the two is one of form, not of
number.

The solver tier was fixed on the published model cards before any
campaign flew, for five reasons, in the order they weighed.

\begin{enumerate}
  \item \textbf{Headroom.} On BrowseComp, the published benchmark closest
        to this bed, Pro scores $83.4$, near the open-model ceiling; a
        solver that strong would flatten the window in which a harness
        effect could show at all. Flash sits mid-curve --- $53.5$ at high
        effort, $73.2$ at maximum --- which leaves an effect somewhere to appear.
  \item \textbf{Cost.} Pro's list price is $3.1\times$ Flash's; after
        Flash's roughly $1.3\times$ longer outputs the per-task ratio is
        near $2.4\times$, which over thousands of rollouts is the
        difference between the replication in \S\ref{sec:replication}
        being affordable and not.
  \item \textbf{Failure signature.} Third-party measurement reported
        Flash's tool calling as clean --- no runaway retries, no malformed
        arguments --- with its failures falling on substantive correctness
        and long-context retrieval. Those are failures a harness can act
        on; protocol failures are not.
  \item \textbf{The one documented hard limit is avoided.} Flash collapses
        on long context with thinking off (MRCR-1M $37.5$); thinking is on
        in every trajectory here, and every trajectory carries a
        non-empty thinking block.
  \item \textbf{Two unknowns were left to the pilot} rather than
        interpolated from cards that do not publish them: strict
        final-answer format compliance, and run-to-run variance at twenty
        steps. The second became the subject of this thesis.
\end{enumerate}

The meta tier was set to Pro so that diagnosis quality would not be the
bottleneck when the treated arm was read: if graph evidence failed to help
a strong reader, the failure could not be blamed on the reader. The
expectation stated before the campaigns was therefore modest and
directional --- a mid-curve solver with room to improve, a strong meta
tier, and an effect, if the evidence mattered, that would show on the
score. Claims are scoped to this pair and stated that way. It is recorded
as a fact about the experiment rather than as a threat to it, because no
claim in this thesis is made about model tiers.

% -------------------------------------------------------------------------
\section{Power arithmetic first}
\label{sec:power}
% -------------------------------------------------------------------------

\subsection*{The flip rate}

Take two consecutive rounds in which nothing shipped. The configuration is
identical; the only difference is that the bed was run again. Across the
six such windows in the baseline campaign's sixteen-round readout,
\textbf{21.0\%} of task outcomes flip --- $126$ of $600$ task-pairs, 95\%
Wilson interval $[17.9\%, 24.4\%]$ \F{15}. The other two seeds give
$20.1\%$ and $19.2\%$, and pooled over all twenty windows the figure is
\textbf{20.1\%} on $402/2{,}000$ pairs, $[18.4\%, 21.9\%]$, with the three
per-seed intervals overlapping heavily \F{9e}. \emph{The noise floor is a
property of the bed, not of a run.}

Over those windows, $52.3\%$ of task-campaign pairs flip at least once.
The resulting total-score swing for a single window runs from $-8$ to
$+5$ tasks \F{15}\F{44}\F{45}.

\subsection*{The arithmetic, walked slowly}

\begin{enumerate}
  \item One fifth of task outcomes flip on re-measurement with no treatment
        applied.
  \item The bed therefore swings by up to eight tasks between two
        identical rounds.
  \item Any total-score delta of five tasks or fewer is consequently
        unreadable as evidence of an effect.
  \item The score axis can therefore return only ``not separable'' for
        per-round effects of the size this loop produces.
\end{enumerate}

Step 4 is the one that matters, and it is not a complaint about sample
size. It is a statement about what this instrument can resolve, and it is
the reason \S\ref{sec:efficacy-protocol} exists.

\subsection*{One envelope, one caliber}

The swing is quoted throughout as an \emph{observed range} over ship-free
windows, and the range depends on how many windows were drawn:
$-4\ldots{+}5$ on seed 1 (6 windows), $-4\ldots{+}2$ on seed 2 (8
windows), and $-8\ldots{+}5$ on seed 3 (6 windows) \F{15}\F{44}\F{45}.
The band used for every readability judgment in this thesis is the union,
$\mathbf{-8\ldots{+}5}$ tasks, quoted that way everywhere; per-seed
ranges appear only where a seed is being described.

\emph{A range over a handful of windows is not a standard deviation, and
the difference is not cosmetic.} With all three seeds on one bed there are
twenty same-config windows, which is enough to estimate the per-round score
SD directly rather than convert a range: it is $\mathbf{3.70}$ tasks, so
three sigma is $\pm 11.1$ --- \emph{wider} than the observed band
\F{15}. Every ``unreadable'' verdict in this thesis is therefore
conservative in the direction that makes separation \emph{easier} to
claim, not harder.
Where the two calibers disagree about a specific number, the disagreement
is stated at that number rather than resolved silently; the one case that
arises is the $+9$ of \S\ref{sec:one-repair}, and it is discussed there and
in \S\ref{sec:readout}.

\subsection*{There is no informative-and-reliable subset of this bed}

The natural response to a noisy bed is to find its reliable core and
measure there. On this bed that response fails, and the census says so
directly. Across the six whitelist campaigns on the 100-task subset ---
9{,}600 task evaluations --- \textbf{3} tasks are never solved,
\textbf{18} always pass, and \textbf{79 are volatile} \F{54}.

The 21 stable tasks contribute no variance and can therefore register no
treatment effect. The 79 that \emph{can} move are precisely the ones that
move unaided. Removing the volatile family does not clean the instrument;
it removes the instrument.

Bed reliability here cannot be bought by item selection. It can be bought
in repetitions --- a Spearman--Brown calculation prices that against the
measured per-task flip rate --- or by changing the readout, which is this
thesis's own route (\S\ref{sec:future}).

\subsection*{Same data, two verdicts}

The design answer is to stop reading totals and start reading pairs. One
window makes the point sharply: a round whose \emph{total} delta was $+7$
--- inside the band, and therefore unreadable --- decomposes into 8 tasks
fixed and 1 broken, which under an exact McNemar test gives $p \approx
0.039$ \F{21}.

Same data, two verdicts. The problem was never too little data. It was the
wrong reading axis.

% -------------------------------------------------------------------------
\section{The efficacy-trial protocol}
\label{sec:efficacy-protocol}
% -------------------------------------------------------------------------

Level 3 of the readout ladder --- did firing change outcomes --- was not
available from the loop, so it was built by hand. The protocol below is
what a manual L3 readout costs, and it is offered as a specification for
the automated version \S\ref{sec:future} calls for.

\begin{itemize}
  \item \textbf{Arms are interleaved within the same window.} Running arms
        serially confounds arm with time window. This rule was learned by
        having a result destroyed by its absence \F{36}.
  \item \textbf{True completions are counted before anything is scored.}
        Launch-spike deaths and phantom runs are excluded by inspection.
        This is registry error 14.
  \item \textbf{Campaign task parameters are replicated exactly.} The
        campaigns run \texttt{--max-steps 20}; the loader's default of 40
        doubles the step budget and revives tasks that had died of burnout,
        which silently changes what is being measured. Registry error 15.
  \item \textbf{Adjudication files are quarantined outside the
        subject-reachable filesystem, and seeded with canary words.} If a
        planted marker word ever surfaces in a subject trajectory, the
        subject read what it should not have. Registry error 13 is the
        specimen: an adjudication file left readable in a run root, a role
        that read the answer sheet, and a round of rulings voided.
  \item \textbf{Gold labels are health-checked before trials}, following
        the archaeology protocol --- a wrong gold label was found this way
        \F{29}.
  \item \textbf{A canary set of eight always-pass tasks detects harm}
        \F{33}, stopping arithmetic decides continuation, and
        \textbf{mechanism endpoints are reported beside scores}: firing
        forensics, and the confident-wrong share --- the fraction of
        failures in which the subject closed confidently on a wrong
        answer.
\end{itemize}

The last item is what makes the protocol able to return an interpretable
null. Scores flat with mechanism flat is a true zero. Scores moving with
mechanism flat is a noise finding. Without the second endpoint the two are
indistinguishable, and \S\ref{sec:clinic} reports one of each.

% -------------------------------------------------------------------------
\section{Resolution laws as design inputs}
\label{sec:resolution}
% -------------------------------------------------------------------------

The bed-level flip rate of \S\ref{sec:power} is the mild face of the
problem. At task level it is worse: per-task base rates swing by up to
$\pm 50$ percentage points between half-hour windows \F{34}. Two
consequences are treated as design constraints rather than as findings:
single-window deltas of five or fewer are unreadable, and any efficacy
claim requires pooling across windows.

\subsection*{Multiplicity, declared}

There is \textbf{exactly one confirmatory test} in this thesis: the
\emph{prospective} efficacy batch of $n{=}70$ pairs --- 7 clean tasks
$\times$ 10 repetitions $\times$ 2 arms --- which is the test the stopping
rule was registered against \F{34}. The pooled 174-pair reading is a
\emph{supporting} estimate over all clean same-window pairs, reported
because it agrees, not because it was registered.

Every other $p$-value in this document is exploratory. That includes the
McNemar window of \S\ref{sec:power} \F{21}, the same-window controls
\F{36}, the notes channel \F{35}, and the cluster transfer \F{32}. Each is
labelled exploratory at its point of use as well as here, and no
family-wise correction is claimed.

This matters because of what the instruments actually did during the study.
Four successive apparent wins --- $+7$, $+5$, $+4$, $+6$ --- were executed
one at a time, by three different instruments. Those four are serial
re-tests of a single hypothesis, which is exactly why the confirmatory
endpoint is the pre-registered prospective batch and not any of them.

% -------------------------------------------------------------------------
\section{Evaluator-bias controls}
\label{sec:evaluator-bias}
% -------------------------------------------------------------------------

\paragraph{pass@1 discipline.} Mixed-$k$ passes do not count. Applying this
rule mechanically to the baseline campaign's seed round moved its recorded
score from 67 to \textbf{57}: $k$-mixing had gifted the seed ten free
passes.

The correction is stated here rather than left for a reader to find,
because of the direction it moves. Lowering the seed \emph{enlarges} the
campaign's apparent gain --- the correction flatters this thesis's own
baseline campaign. The rule is applied mechanically in both directions, and
this instance is disclosed precisely because it is the instance where the
mechanical rule happens to help.

\paragraph{The error registry.} Fifteen errors of measurement and design
were logged during this work, with forced re-checks after each. Both bias
directions are documented. Three of the fifteen --- errors 13, 14 and 15,
all of which appear in \S\ref{sec:efficacy-protocol} --- were caught not by
the author but by the supervisor, and the root cause common to them is
recorded as \emph{stopping the check too early}.

The registry is treated as a result in its own right and is discussed in
\S\ref{sec:threats}. A study whose evaluator is not itself audited has no
grounds for the confidence intervals it reports.
